% =====================================================================
\chapter{Conclusions and Future Work}
\label{ch:conclusions}
% =====================================================================

\section{Limitations}
\label{sec:concl-limitations}

This study is deliberately scoped, and several limitations should be read
alongside its findings. First, the controlled diagnosis centres on a single
primary backbone, LayoutLMv3-base; although LiLT and BROS are used to test
generalisation on a reduced subset, the per-component picture is established most
firmly for one architecture. Second, the task sequence comprises three public
datasets (FUNSD, CORD, SROIE); while they span forms and receipts and a wide range
of label-space sizes, they are modest in scale and do not cover every document
genre. Third, the text-only diagnostic condition (C1) is a same-architecture
LayoutLMv3 with its non-text pathways zeroed rather than a true unimodal encoder,
so it is the closest apples-to-apples text baseline rather than an independent
one. Fourth, the work targets token-level key-information extraction and does not
address entity linking, relation extraction, or generative document IE. Fifth, all
results use the OCR transcriptions and bounding boxes supplied with the benchmark
datasets and never re-run a recogniser; this is a deliberate choice that removes OCR
quality as a confounder for the diagnosis, but it also means the findings are
established under clean, dataset-provided text and layout. A deployed pipeline would
introduce its own OCR engine, whose recognition and layout-detection errors differ
across document genres and could interact with the per-component forgetting picture;
quantifying that interaction is left to future work. Finally, results rest on three
seeds; more seeds would tighten the confidence intervals around the smaller
effects.

\section{Conclusions}
\label{sec:concl-conclusions}

This thesis set out to answer \emph{where} a tri-modal document encoder forgets,
and to use that answer to design a continual-learning remedy. Its first
contribution is a per-component forgetting diagnostic that attributes catastrophic
forgetting to the text, visual, layout, and fusion components of LayoutLMv3 under
a controlled modality-ablation protocol, using complementary representational
(CKA) and importance-based (Fisher) metrics together with standard forgetting
measures.

The diagnosis returns a clear and somewhat counter-intuitive finding: under naive
sequential fine-tuning LayoutLMv3 forgets almost totally---an earlier task that
scores $88$ entity-F1 collapses to $\sim$2 F1 after a single subsequent task---and
this forgetting is \emph{output-side}, not fusion-specific. It concentrates in the
classifier head and the late encoder layers: the head dominates per-group Fisher
importance by roughly an order of magnitude, and per-layer CKA drift rises
monotonically with depth (from $1.00$ at the input encoders to $0.22$ at the head),
while the input encoders are essentially unchanged. The cross-condition test could
not statistically distinguish the full multimodal model's forgetting from its
unimodal ablations, so the original premise that multimodal encoders forget
differently \emph{in fusion} is not supported; the locus is the output head.

The benchmark of continual-learning strategies confirms what this diagnosis predicts.
Across the three scenarios the ordering is unambiguous---replay $\gg$ regularisation
$\gg$ distillation---and it follows directly from \emph{where} forgetting lives:
because the damage is at the head, \emph{replay}, which re-exposes old-class targets at
the output, is the most effective remedy. On domain-incremental learning both replay
methods (ER, DER++) reach the joint upper bound (AA $\approx 88$ vs the oracle's $85$,
backward transfer $\approx 0$), closing essentially the entire $\sim$47-point forgetting
gap; EWC roughly halves forgetting wherever a domain shift is present (DIL/Mixed BWT
$-37$/$-31$ vs naive $-71$/$-67$); and LwF tracks the naive lower bound throughout, as
its distilled old-class logits decay once the growing head stops reinforcing them.

The proposed method, DocCL---a mechanism that would concentrate continual-learning
effort on the dominant-forgetting head and late layers---was designed from the
diagnosis but \emph{not} empirically evaluated in this work: the diagnosis was
inconclusive as to which single architectural lever dominates, so per a pre-registered
fallback the diagnosis and benchmark are reported as the primary contribution and the
full DocCL evaluation is left as the immediate next step (\S\ref{sec:concl-future}).
The diagnostic contribution stands on its own: it provides a reusable tool and a shared
benchmark for studying forgetting in multimodal document encoders---together with a
strong-baseline context (replay reaching the oracle) that any future targeted method
must beat---which the prior, fragmented literature lacked.

\section{Future Work}
\label{sec:concl-future}

Several extensions follow naturally and are reserved for post-thesis work; each is
a future direction rather than part of the present AAAI-phase scope.

\emph{Scaling and broadening the diagnosis.} The controlled diagnosis is
established for an encoder-only, base-scale backbone. A natural next step is to
re-run the same per-component protocol on larger and architecturally different
backbones---LayoutLMv3-large, the multilingual LayoutXLM, and decoder-based or
instruction-tuned document MLLMs such as InternVL3---to test whether the
dominant-forgetting component is a property of the LayoutLMv3-base design or a more
general feature of tri-modal document encoders. Decoder-based encoders are
especially interesting because their generative objective and tokenisation differ
markedly from the sequence-labelling setup used here.

\emph{Additional continual scenarios.} The three studied scenarios cover
class-incremental and domain-incremental shifts. Future work could add
task-incremental learning, where a task identifier is available at inference, and
combined class-and-domain shifts that vary the label space and the document genre
simultaneously; both would stress the components differently and test whether a
single mechanism-targeted remedy transfers across regimes.

\emph{Cross-lingual continual document IE.} Extending the benchmark with a
multilingual corpus such as XFUND would layer language shift on top of domain
shift, asking whether forgetting concentrates in the text pathway when the script
and vocabulary change while layout statistics stay comparatively stable---a setting
the current English-only benchmark cannot probe.

\emph{Richer, layout-aware remedies.} The diagnosis points to a specific component;
a complementary direction is to enrich the remedy on the data side with a
document-aware structural replay that rehearses not raw tokens but layout-structured
exemplars (for example region- or block-level representations), exploiting the very
spatial structure that distinguishes document IE from plain-text IE. Such a replay
could be combined with the mechanism-targeted approach rather than replacing it.

\emph{Joint multi-task continual IE.} Finally, learning token-level extraction
together with relation extraction (and, eventually, entity linking) in one
continual framework is an open and high-value direction: it would move the
benchmark from single-task sequence labelling towards the joint document
understanding that downstream applications ultimately require.

These extensions are intentionally outside the present scope, which prioritises a
rigorous diagnosis and a single, well-justified remedy over breadth; they sketch
the trajectory for a fuller follow-up study.
